% proposicion_brecha_convencion.tex -- Proposición de brecha de convención (A-M22, primer resultado propio)
% Borrador v0, 2026-09-22 (sesión con Francisco 19:00 Praga). Compila solo: pdflatex proposicion_brecha_convencion.tex
% Etiquetas §25 del plan maestro en cada enunciado. Toda cifra citada está en el pie "Números verificados".
% Nota: *.tex está ignorado globalmente (.gitignore:74); para versionar: git add -f augmented/paper/proposicion_brecha_convencion.tex
\documentclass[11pt]{article}
\usepackage[utf8]{inputenc}
\usepackage[T1]{fontenc}
\usepackage{amsmath,amssymb,amsthm}
\usepackage{booktabs}
\usepackage[margin=1in]{geometry}
\usepackage{xcolor}
\newtheorem{proposition}{Proposition}
\newtheorem{corollary}{Corollary}
\newtheorem{lemma}{Lemma}
\theoremstyle{remark}
\newtheorem{remark}{Remark}
\newcommand{\pz}{\mathrm{pz}}
\newcommand{\str}{\mathrm{str}}
\newcommand{\OPT}{\mathrm{OPT}}
\newcommand{\lab}[1]{\textcolor{blue!60!black}{\textsf{\small [#1]}}}
\newcommand{\Cstr}{C_{\str}}
\newcommand{\Cpz}{C_{\pz}}

\title{The convention gap: strict versus posterior-zero clearing\\
\large Proposición de brecha de convención --- borrador v0}
\author{A (statement and proofs of Parts 1 and 3) \and formal version refined with LLM support}
\date{2026-09-22}

\begin{document}
\maketitle

\section{Setting and the two conventions}

Augmented laminar model with a hard budget of $B$ tests on every path. Person $i$ is healthy with probability $q_i$ and infected with $p_i = 1-q_i$, independently; utility $u_i$ is collected once, when $i$ is accredited healthy. A test on a pool $T$ ($|T|\le G$) returns $R(T)$, the exact number of infected in $T$. A \emph{history} $h$ is the sequence of pairs $(T, R(T))$ observed so far. A policy $\pi$ is a decision tree over histories; laminar legality (a new pool is disjoint from every tested pool, or a strict subset of a residual atom) depends only on $h$, never on the accreditation rule.

For a history $h$ write $C(h)$ for the \emph{accredited set} under the rule in force:
\begin{itemize}
  \item \textbf{Strict (hard clearing).} $i \in \Cstr(h)$ iff some tested pool $T$ with $i\in T$ has $R(T)=0$.
  \item \textbf{Posterior-zero (soft clearing).} $i \in \Cpz(h)$ iff $\Pr(i \text{ infected}\mid h)=0$, i.e.\ no infection profile consistent with $h$ has $i$ infected; deduction by subtraction of counts is included.
\end{itemize}
The value of a policy is its expected collected utility, $V_\bullet(\pi)=\mathbb E\big[\sum_{i\in C_\bullet(h_{\rm final})} u_i\big]$, and $\OPT_\bullet(B)=\max_\pi V_\bullet(\pi)$, for $\bullet\in\{\str,\pz\}$. Posterior-zero is the normative convention of the paper (gate G0, 2026-08-25); strict is the comparison variant.

\section{The proposition}

\begin{proposition}[Convention gap]\label{prop:gap}
\leavevmode
\begin{enumerate}
\item[\textbf{1.}] \textbf{Domination.} \lab{DEMOSTRADO} For every history $h$, $\Cstr(h)\subseteq\Cpz(h)$. Hence $V_\str(\pi)\le V_\pz(\pi)$ for every policy $\pi$, and $\OPT_\str(B)\le\OPT_\pz(B)$ for every $B$.

\item[\textbf{2.}] \textbf{Simulation with overhead.} \lab{ENUNCIADO; prueba pendiente} Let $D(h)$ be the set accredited by deduction only, $D(h)=\Cpz(h)\setminus\Cstr(h)$, and let $\Delta(\pi)$ be the maximum over trajectories of $\pi$ of the number of tests needed to certify $D(h_{\rm final})$ under the strict rule (one test per deduced-clean residual piece of size at most $G$). Then every posterior-zero policy is replicated under strict clearing with $\Delta(\pi)$ additional tests, so
\[
\OPT_\pz(B)\;\le\;\OPT_\str\big(B+\Delta\big),\qquad \Delta=\max_{\pi}\Delta(\pi).
\]
Witness in the four-person game of Part 3(b): $\OPT_\pz(2)=0.774\le\OPT_\str(3)=1.011$. \lab{VERIFICADO}

\item[\textbf{3.}] \textbf{Sharpness.}
\begin{enumerate}
\item[(a)] \emph{State witness.} \lab{DEMOSTRADO} A pair with count $1$ and one test left, $u\equiv1$: refining one member is worth $1$ under posterior-zero (whichever way the count falls, someone is proven healthy) and $\tfrac12$ under strict.
\item[(b)] \emph{Minimal full game.} \lab{DEMOSTRADO $B=2$; VERIFICADO $n=4$ por dos enumeraciones independientes} Let $n\ge3$, $G=2$, $B=2$, $u\equiv1$, homogeneous $q$, $p=1-q$. Let $\pi_{\rm par}$ test the pair $\{1,2\}$, refine $\{1\}$ on count $1$, and test the fresh singleton $\{3\}$ otherwise; let $\pi_{\rm sing}$ test two singletons. Then
\[
V_\pz(\pi_{\rm par}) = 2q^3-2q^2+3q = 2q + q\,(q^2+p^2),\qquad
V_\str(\pi_{\rm par}) = 2q^3-q^2+2q = 2q + q^2(2q-1),\qquad
V(\pi_{\rm sing}) = 2q,
\]
so that
\[
V_\pz(\pi_{\rm par})-V_\str(\pi_{\rm par}) = pq,\qquad
V_\pz(\pi_{\rm par})-V(\pi_{\rm sing}) = q\,(q^2+p^2)>0\ \ \forall q\in(0,1),\qquad
V_\str(\pi_{\rm par})-V(\pi_{\rm sing}) = q^2(2q-1).
\]
The first difference comes from the single branch $R=1$. For $q\le\tfrac12$, $\pi_{\rm par}$ is the best pair-first continuation under both conventions and $\pi_{\rm sing}$ is the strict optimum; for $q>\tfrac12$ the two formulas are strict lower bounds (Remark~\ref{rem:region}).
\item[(c)] \emph{Scale.} \lab{cota realizable} At the anchor $(q,G,B)=(0.05,16,7)$ the fixed cover-then-bisect plan is worth $1-0.95^{48}=0.9147$ under posterior-zero and $1-0.95^{32}=0.8063$ under strict: the strict rule spends one full pool of tests on certification. This plan is \emph{not} optimal (Remark~\ref{rem:anchor}).
\end{enumerate}
\end{enumerate}
\end{proposition}

\begin{corollary}[The threshold of the non-augmented model disappears]\label{cor:threshold}
In the two-test game of Part 3(b), strict clearing reproduces the threshold of Proposition~1 of the non-augmented model (arXiv:2206.10660 v4; companion Thm~7.1): for $q<\tfrac12$ the strict optimum never groups, at $q=\tfrac12$ it ties, and for $q>\tfrac12$ it groups. Under posterior-zero there is no threshold: the pair strictly beats the singletons for every $q\in(0,1)$. The convention flips the optimal first move, not only the value. \lab{DEMOSTRADO $B=2$; VERIFICADO $n=4$}
\end{corollary}

\section{Proofs}

\begin{proof}[Proof of Part 1]
If $i\in T$ and $R(T)=0$ then no member of $T$ is infected, so every profile consistent with $h$ has $i$ healthy and $\Pr(i\text{ infected}\mid h)=0$: $\Cstr(h)\subseteq\Cpz(h)$. A policy generates the same histories with the same probabilities under both rules, because legality and transitions depend only on $h$; its value is the expected utility of the accredited set, which is pointwise larger under posterior-zero. Taking the maximum over policies gives $\OPT_\str(B)\le\OPT_\pz(B)$.
\end{proof}

\begin{proof}[Mechanism of Part 2 (proof pending)]
Run the posterior-zero policy $\pi$ under the strict rule: every action legal under posterior-zero is legal under strict (the strict menu additionally contains tests on deduced-clean pieces), so histories and probabilities coincide. At the end of each trajectory, test each deduced-clean residual piece: such a test is legal, reads $0$ with certainty, and accredits the piece under the strict rule; pieces belonging to different atoms cannot be merged into one test (laminarity), which is why the overhead counts pieces. The budget is hard on every path, so the overhead is the \emph{maximum} over trajectories, not the expectation. On the trajectory ``pair reads 1, then $\{1\}$ reads 1'' the extra test is $\{2\}$; it costs one test and pays person~2. \textit{Pending: the exact bound on $\Delta$ (one extra test per refinement event gives $\Delta\le B-1$) and the formal write-up.}
\end{proof}

\begin{proof}[Proof of Part 3(b)]
Branch probabilities of the first test $\{1,2\}$: $\Pr(R=0)=q^2$, $\Pr(R=1)=2pq$, $\Pr(R=2)=p^2$. Branch totals of $\pi_{\rm par}$ (already accredited plus expected gain of the second test):
\begin{center}
\begin{tabular}{lccc}
\toprule
branch & already accredited & second test & branch total \\
\midrule
$R=0$ & $2$ & fresh singleton $\{3\}$: $q$ & $2+q$ \\
$R=1$ & $0$ & refine $\{1\}$: $1$ (pz) / $\tfrac12$ (str) & $1$ / $\tfrac12$ \\
$R=2$ & $0$ & fresh singleton $\{3\}$: $q$ & $q$ \\
\bottomrule
\end{tabular}
\end{center}
In branch $R=1$, under posterior-zero the refinement pays in both outcomes (count $0$: person 1 healthy; count $1$: person 2 healthy by deduction), under strict only in the first. Weighting by the branch probabilities gives the three polynomials; the differences follow by expanding with $p=1-q$, using $2q^2-2q+1=q^2+p^2$. The fresh singleton is the best second test in the resolved branches when $q\ge 2q^2$, i.e.\ $q\le\tfrac12$ (the fresh pair $\{3,4\}$ is worth $2q^2$: it pays only when both are healthy, since a count of $1$ with no budget left accredits nobody). Under strict, in branch $R=1$ the refinement ($\tfrac12$) beats the fresh singleton ($q$) exactly when $q\le\tfrac12$. The strict optimality of $\pi_{\rm sing}$ for $q\le\tfrac12$ was checked by exhaustive enumeration at $n=4$ (see below).
\end{proof}

\begin{remark}[Region of exactness and lower bounds]\label{rem:region}
For $q>\tfrac12$ the best second test in the resolved branches is the fresh pair ($2q^2>q$), and for $q>1/\sqrt2$ also in branch $R=1$ ($2q^2>1$); the polynomials of Part 3(b) are then the value of a legal policy, hence strict lower bounds on the best pair-first value (e.g.\ $q=0.6$: formula $1.512$, exact pair-first value $1.5744$). The claim of the corollary lives in $q\le\tfrac12$, the region of Nick's Proposition~1, where the formulas are exact.
\end{remark}

\begin{remark}[Why $n\ge3$]
With $n=2$ the second test has nowhere to go in the resolved branches: $V_\pz(\pi_{\rm par})=2q^2+2pq=2q=V(\pi_{\rm sing})$, an exact tie, while $V_\str(\pi_{\rm par})=2q^2+pq=q(1+q)<2q$. The deduction alone makes the pair worth $2q$; the third person, who absorbs the second test in the resolved branches (probability $q^2+p^2$, gain $q$), turns the tie into the strict gain $q(q^2+p^2)$.
\end{remark}

\begin{remark}[Three tests]
The threshold of the corollary is a $B=2$ statement. With $B=3$ at $q=0.3$: $\OPT_\str(3)=1.011>0.9$, so strict clearing already groups; and under posterior-zero the pair-first and singleton-first optima tie exactly at $537/500=1.074$. The advantage of grouping is not monotone in $B$.
\end{remark}

\begin{remark}[The anchor plan is not optimal]\label{rem:anchor}
Map 3 (B, 2026-09-21, \texttt{results/mapa\_homogeneo\_3.csv}): at the frontier $n=24$, $G=8$, $B=6$ the exact laminar optimum is $0.8654$ against the cover-then-bisect bound $0.7080$; in the anchor population with actions restricted to $G\le8$ the optimum is at least $1.1233$ against the bound $0.9147$; the full anchor $G=16$ is not certified (state limit). Part 3(c) is therefore a statement about the fixed plan under the two conventions, not about optima; the anchor bound is slack.
\end{remark}

\section{Verification by enumeration}

Two programs that share no code compute the exact laminar optimum: the Bellman solver on residual atoms (\texttt{augmented/bm17\_toy\_solver.py}) and brute-force enumeration over explicit profiles and histories (\texttt{augmented/enumerador\_historias.py}; independently, B's \texttt{augmented/pathwise\_enum.py}). They agree exactly, in rational arithmetic, on every optimum and every forced first action for $n\le5$, $B\le3$, both conventions (\texttt{tests\_brecha\_convencion.py}, \texttt{tests\_bm17.py}: 626 tests). The table gives $n=4$, $G=2$, $B=2$, $u\equiv1$, where the formulas of Part 3(b) are exact:
\begin{center}
\begin{tabular}{cccccc}
\toprule
$q$ & $V_\pz(\pi_{\rm par})$ & $V_\str(\pi_{\rm par})$ & $V(\pi_{\rm sing})=2q$ & best first move, pz & best first move, str \\
\midrule
0.1 & 0.282 & 0.192 & 0.2 & pair & singleton \\
0.2 & 0.536 & 0.376 & 0.4 & pair & singleton \\
0.3 & 0.774 & 0.564 & 0.6 & pair & singleton \\
0.4 & 1.008 & 0.768 & 0.8 & pair & singleton \\
0.5 & 1.250 & 1.000 & 1.0 & pair & tie \\
\bottomrule
\end{tabular}
\end{center}
At scale, Map 1 (B, 2026-09-21, \texttt{results/mapa\_homogeneo\_1.csv}): for $q\in\{0.05,\dots,0.95\}$, $G\in\{2,3,4,5,8\}$, $B\in\{2,\dots,8\}$, $n=BG$, $u\equiv1$, posterior-zero, in $665$ of $665$ cells the optimal first action is a group. \lab{VERIFICADO por B}

\section{A's version, in his own words}

\begin{quote}\small
\begin{enumerate}
\item We compare two rules for collecting utility. Soft clearing accredits a person as healthy as soon as the whole history proves it, whatever the proof. Hard clearing accredits a person only if they sat in a tested pool whose count came back zero.
\item Soft clearing never collects less than hard clearing: a zero count on your pool is one of the proofs soft clearing accepts, so every hard-clearing certificate is also a soft-clearing certificate. With a single test the two rules collect the same; the difference needs a deduction.
\item Soft clearing is not cheating. Everything it collects by deduction, a hard-clearing player could also collect by spending extra tests on people already known to be healthy, tests whose result is known before running them. So the soft-clearing optimum with budget $B$ never exceeds the hard-clearing optimum with a somewhat larger budget; in the four-person game one extra test is enough, 0.774 with two tests against 1.011 with three.
\item The gap is real and it comes from one event. Test a pair, get count one, then test one member: soft clearing pays with certainty, because whichever way the count falls someone is proven healthy; hard clearing pays half the time. With four people, health probability 0.3 and two tests, the same policy is worth 0.774 under soft clearing and 0.564 under hard clearing, while two individual tests are worth 0.6.
\item So the convention flips the optimal first move, not only the value: under hard clearing you should never group, under soft clearing you should open the pair.
\item In the two-test game hard clearing reproduces the threshold of Nick's Proposition~1: below one half the optimum never groups, above it groups. Under soft clearing there is no threshold, the pair beats the individuals at every probability.
\item At scale Héctor's Map 1 says the same: in 665 of 665 homogeneous cells the optimal first move under soft clearing is a group.
\end{enumerate}
\end{quote}

\medskip\noindent\textit{Números verificados en este documento: $3/10$, $3/5$, $387/500=0.774$, $141/250=0.564$, $1011/1000$, $537/500$ (referencias de diseño B-M17, dos vías); tabla $q\in\{0.1,\dots,0.5\}$ y valor exacto $1.5744$ en $q=0.6$ (enumeración, 2026-09-13; \texttt{tests\_brecha\_convencion.py} fija las fórmulas para $q\le\tfrac12$ y las desigualdades para $q>\tfrac12$); $1-0.95^{48}=0.9147$, $1-0.95^{32}=0.8063$ (aritmética, §16 del plan); Mapa 1 $665/665$ y Mapa 3 $0.8654$/$0.7080$/$1.1233$ (B, commit 5719df3). Pendiente: prueba de la Parte 2.}

\end{document}
